\paragraph{Justificación Metodológica de la Selección de Module-LWE para el Prototipo}

La elección del problema matemático base para el desarrollo del prototipo de autenticación en Python debe equilibrar estrictamente tres criterios: seguridad teórica, dimensiones de memoria y velocidad computacional \cite{yang2025multiserver, mukil2026implementing, ma2026laasso}. A continuación se justifican las razones por las cuales Module-LWE constituye la mejor alternativa de diseño frente a sus competidoras:

\begin{itemize}
    \item \textbf{Superación del Sobrecosto de Memoria de LWE Estándar:} En LWE puro, para alcanzar niveles de seguridad poscuántica adecuados ($n \ge 768$), la matriz pública $\mathbf{A} \in \mathbb{Z}_q^{n \times m}$ requiere almacenar cientos de miles de enteros modulares, saturando la memoria del cliente y del servidor \cite{yang2025multiserver}. En MLWE, al agrupar los coeficientes en polinomios de grado $d=256$, una pequeña matriz polinomial de dimensión $k \times k = 3 \times 3$ logra el mismo grado de expansión algebraica ($n = k \cdot d = 768$) reduciendo el almacenamiento de la clave pública a aproximadamente $1.15$~KB \cite{yang2025multiserver}.
    
    \item \textbf{Mitigación de Vulnerabilidades Algebraicas de RLWE:} Si bien RLWE ($k=1$) ofrece la máxima compresión al operar sobre un único anillo polinomial, su estructura de retículo ideal puro introduce simetrías algebraicas adicionales que podrían ser potencialmente explotables por algoritmos de reducción geométricos específicos \cite{yang2025multiserver}. MLWE rompe esta dependencia de ideales puros mediante la introducción de la dimensión matricial $k \ge 2$, lo que blinda la estructura contra vulnerabilidades de retículos ideales sin sacrificar el rendimiento \cite{yang2025multiserver, ma2026laasso}.
    
    \item \textbf{Eficiencia Algorítmica en Python:} MLWE preserva intacta la capacidad de acelerar el producto de matrices de polinomios $\mathbf{A}\mathbf{s}$ mediante la Transformada de Teoría de Números (NTT) en tiempo $O(d \log d)$ por cada entrada de la matriz \cite{nguyen2019high, yang2025multiserver}. Esto permite que las operaciones de encriptación y prueba de conocimiento cero (ZKP) en el prototipo en Python se ejecuten en pocos milisegundos \cite{mukil2026implementing, ma2026laasso}, garantizando la viabilidad práctica y la solidez poscuántica estandarizada por el NIST \cite{ma2026laasso}.
\end{itemize}
